% ChargeShield-FL --- DSN 2027 --- 2. Background and Related Work
% Scritta il 2026-09-23. Assorbe la parte ancora valida di preliminaries.tex
% e related_work.tex, con tre correzioni: il protocollo usa FedProx con
% mu = 0.01 (related_work.tex diceva mu = 0); il campo operativo e' T = 10
% round (preliminaries.tex usava T = 3, che e' il regime canary); esiste
% anche il meccanismo record-level DP-SGD. Filoni e implicazioni dalla guida,
% sez. 3.

\section{Background and Related Work}
\label{sec:background}

\subsection{Differential Privacy and the Protected Unit}
\label{sec:bg-dp}

A randomised mechanism $\mathcal{M}$ satisfies $(\varepsilon,\delta)$-DP if,
for every pair of adjacent datasets $D, D'$ and every set of outputs $O$,
\begin{equation}
\Pr[\mathcal{M}(D) \in O] \le e^{\varepsilon} \Pr[\mathcal{M}(D') \in O] + \delta
\label{eq:dp}
\end{equation}
\cite{dwork2006calibrating,dwork2014algorithmic}. The definition is empty until
the adjacency relation is fixed, and the adjacency is what we call the
\emph{protected unit}. Under \emph{record-level} adjacency $D$ and $D'$ differ
in one record; DP-SGD achieves it by clipping each example's gradient, adding
Gaussian noise to the sum, and tracking the privacy loss across steps with a
moments or R\'enyi accountant~\cite{abadi2016deep,mironov2017renyi}. Under
\emph{client-level} adjacency they differ in one client's data; DP-FedAvg
achieves it by clipping each client's update and adding noise to the
aggregate~\cite{mcmahan2018learning}. Membership inference, as we measure it,
concerns records. The two mechanisms therefore relate differently to the
quantity we measure, and the same $\varepsilon$ does not describe equivalent
protection under the two (Section~\ref{sec:system-dp}).

Composition bounds the guarantee of $T$ rounds: basic composition gives
$(T\varepsilon, T\delta)$, and advanced composition a bound growing as
$\sqrt{T}$ that is tighter only for small $\varepsilon$ and large
$T$~\cite{dwork2014algorithmic}. Under the hypothesis-testing reading of DP,
an $(\varepsilon,\delta)$ guarantee bounds any membership test on the
protected unit: $\mathrm{TPR} \le e^{\varepsilon}\mathrm{FPR} + \delta$, hence
an advantage $\mathrm{TPR} - \mathrm{FPR}$ of at most
$(e^{\varepsilon} - 1 + 2\delta)/(e^{\varepsilon} + 1)$~\cite{kairouz2015composition}.
The bound assumes that the protected units are independent; records of the
same person are not, and the bound need not hold for
them~\cite{humphries2020differentially}.

\subsection{Membership Inference and Its Metrics}
\label{sec:bg-mia}

Given a target record $x$ and a model $f$, an attack computes a score $s(x,f)$
meant to be larger for members than for non-members. For an autoencoder the
natural per-record quantity is the reconstruction error
$\ell(x,f) = \lVert f(x) - x \rVert_2^2 / d$. The \emph{loss-threshold} attack
uses $s = -\ell$~\cite{yeom2018privacy}; a \emph{shadow} attack calibrates the
threshold on models trained by the attacker on data of known
membership~\cite{shokri2017membership,salem2019mlleaks}; LiRA fits, for each
record, Gaussian distributions of its loss under shadow models trained with it
(IN) and without it (OUT) and scores the log-likelihood
ratio~\cite{carlini2022membership}. Sablayrolles et al.\ showed that, when the
two loss distributions are known, a function of the loss is Bayes-optimal, so
access beyond the loss helps only by improving their
estimate~\cite{sablayrolles2019whitebox}.

We report the AUC for continuity with the literature and the true-positive
rate at a fixed false-positive rate (TPR at FPR), whose chance level is
$\mathrm{TPR} = \mathrm{FPR}$, not $0.5$. A near-chance AUC does not exclude a
residual in the low-FPR region of the ROC or in a subgroup of
records~\cite{carlini2022membership}, which is why we also analyse records
individually (Section~\ref{sec:method-perrecord}).

\subsection{Related Work}
\label{sec:related}

\paragraph{DP, attacks and utility} The empirical trade-off between
protection, membership inference and model quality is well studied. Jayaraman
and Evans~\cite{jayaraman2019evaluating} showed that the $\varepsilon$ values
used in practice are far from those at which attacks degrade, and that
evaluating a DP model requires reporting both attack success and utility.
Rahman et al.~\cite{rahman2018membership} measured membership inference
against differentially private deep models, and Naseri
et al.~\cite{naseri2022local} compared local and central DP in FL against
membership inference and backdoor attacks. We do not claim the first DP/MIA
study, in FL or elsewhere; what we add is a controlled comparison in one real
system, across protected units and observation points, with the cost on the
model reported alongside.

\paragraph{Attack quality and metrics} LiRA and the evaluation at low FPR
address the limits of average metrics~\cite{carlini2022membership}. We adopt
both, and we report where the per-record calibration fails in our regime
(Section~\ref{sec:results-canary}).

\paragraph{Validity of the evaluation} Several works document apparent
membership signals produced by differences between the member and non-member
distributions rather than by training. Duan et al.~\cite{duan2024mia} found
that attacks on large language models are close to chance once members and
non-members are drawn from the same distribution; Gomez~\cite{gomez2026clip}
showed that attacks on CLIP relied on out-of-distribution non-members and
collapse to chance on a controlled benchmark; Bai et al.~\cite{bai2026splitbias}
showed that audits of synthetic medical images can detect the data split
rather than membership. The validity of the measurement is therefore a known
problem. We draw members and non-members from the same site pools with the
same split procedure, and we do not present that correction as a new method.

\paragraph{Canaries and auditing} Canaries and adversary instantiation turn a
DP mechanism into an empirical lower bound on $\varepsilon$ under stated
assumptions~\cite{jagielski2020auditing,nasr2021adversary,nasr2023tight};
Steinke et al.~\cite{steinke2023privacy} obtained such bounds from a single
training run with many randomly included canaries; CANIFE crafts canaries for
FL~\cite{maddock2023canife}; Andrew et al.~\cite{andrew2024oneshot} estimate
the empirical privacy of federated training in one run with random canary
clients. Our canary is simpler and serves a different purpose: it tests
whether our attacks detect induced memorisation, and it yields no bound on
$\varepsilon$.

\paragraph{Subject-level membership} Suri et al.~\cite{suri2022subject}
studied whether a subject whose records are spread across silos took part in
federated training. Users who charge at several sites exist in our data
(Section~\ref{sec:system-data}); that motivates a limit of our study, not a
subject-level attack we evaluate.

\paragraph{EV charging} PFAD-BC~\cite{aldweesh2026pfadbc} trains a federated
LSTM autoencoder for anomaly detection on ACN-Data and on a labelled
EV-charging attack dataset, with DP and Byzantine-resilient aggregation, and
evaluates membership inference with a metric that is not an AUC. It is the
closest prior work in domain; we compare protocols with it, not numbers.

\paragraph{Our prior work} Our QRS 2026 fast abstract~\cite{imperatrice2026toward}
stated the problem of the information exposed by FL artifacts in OT
infrastructures. The present paper measures it.

Table~\ref{tab:comparability} places the closest works along the dimensions
that determine whether their results transfer to ours.

\begin{table*}[t]
\centering
\caption{Comparability with the closest prior work. n.e.: not extracted in
this draft; n.a.: not applicable.}
\label{tab:comparability}
\footnotesize
\begin{tabular}{@{}p{2.6cm}p{2.0cm}p{2.3cm}p{3.1cm}p{2.6cm}p{3.3cm}@{}}
\toprule
Work & Setting & Protected unit & Adversary's observation & Controls / audit & Difference from this study \\
\midrule
Jayaraman and Evans~\cite{jayaraman2019evaluating} & centralised & record (DP-SGD variants) & model outputs and loss & n.e. & centralised classifiers; no federated observation points \\
Naseri et al.~\cite{naseri2022local} & FL & client (local and central DP) & n.e. & n.e. & no reconstruction cost on real OT data \\
Carlini et al.~\cite{carlini2022membership} & centralised & n.a. (attack paper) & outputs and loss, shadow models & n.a. & defines LiRA and low-FPR evaluation, which we adopt \\
CANIFE~\cite{maddock2023canife} & FL & client & client updates & crafted canaries, $\varepsilon$ estimate & our canary tests sensitivity, not $\varepsilon$ \\
Andrew et al.~\cite{andrew2024oneshot} & FL & client & n.e. & random canary clients, one run & no natural-regime attack comparison \\
Suri et al.~\cite{suri2022subject} & cross-silo FL & subject & n.e. & n.e. & we measure multi-site users but do not attack subjects \\
PFAD-BC~\cite{aldweesh2026pfadbc} & FL, ACN-Data & n.e. & n.e. & n.e. & different model, metric and aggregation \\
This work & cross-silo FL, 3 real sites & client and record & global model; per-client update at A0--A3 & canary, role exchange, untrained baseline & --- \\
\bottomrule
\end{tabular}
\end{table*}

\Susy{Matrice di confrontabilita' da completare dal testo originale dei paper
(guida sez. 3: task/modello, split, unita', osservazioni, algoritmo, budget e
accounting, metriche, controlli, qualita' del modello, codice). Le celle n.e.
non sono state estratte: non riempirle per analogia.}
